\section{Consequences for generation and free entropy}
\label{sec:consequences}

The universal statement of Theorem~\ref{thm:main} has consequences
beyond the relative generation theorem. We first apply classical
reductions, then compare two generating tuples of a free group factor.

\subsection{General von Neumann algebras and generator invariants}

Willig's direct-integral reduction~\cite[Corollary~2]{Willig1974},
together with the established type~I and properly infinite cases,
reduces the generator problem for von Neumann algebras with separable
predual to the $\mathrm{II}_1$ factor case. Theorem~\ref{thm:main}
therefore gives single generation for every von Neumann algebra with
separable predual.

Let $G$ and $G_{\mathrm{sa}}$ be the generator invariants defined in
\cite[Definition~2.7]{DSSW2008}. They are nonnegative infima over
finite generating sets and finite projection partitions;
$G_{\mathrm{sa}}$ restricts the generating sets to self-adjoint
elements. For every complex $\mathrm{II}_1$ factor $M$ with separable
predual and normalized trace, we have
\[
 G(M)=G_{\mathrm{sa}}(M)=0.
\]
Indeed, by \cite[Corollary~5.1]{DSSW2008}, a positive $G$-value for
one such factor would force a nonsingly generated factor in the same
class, contrary to the universal assertion of
Theorem~\ref{thm:main}. The cited generation convention need not
include the unit. To pass to that convention, replace a generator
$x$ by $y=x+\lambda1$ with $\lambda>\norm{x}$. The invertible $y$
has support projection $1$ in its possibly nonunital von Neumann
closure; that closure then contains $x$ and equals $M$.
Nonnegativity gives $G(M)=0$, and
\cite[Theorem~5.5]{DSSW2008} gives
$G_{\mathrm{sa}}(M)=2G(M)=0$. These values are the defining infima;
no generating set or projection partition attaining zero is asserted.

\subsection{Dependence of free entropy dimensions on generators}

For a bounded finite self-adjoint tuple $Z$ in a tracial von Neumann
algebra, let $\delta(Z)$ and $\delta_0(Z)$ denote Voiculescu's
microstates free entropy dimension and its modified version,
respectively~\cite[Section~2.4]{Voiculescu2002Entropy}.
The $W^*$-invariance question asks whether either number is the same
for any two finite self-adjoint tuples generating the same tracial
von Neumann algebra; see~\cite[Section~2.6]{Voiculescu2002Entropy}.

\begin{corollary}\label{cor:entropy-generators}
For every integer $n\geq3$, the factor $L(\mathbb F_n)$ with its
canonical trace has two finite self-adjoint generating tuples on
which both $\delta$ and $\delta_0$ take different values. In
particular, neither dimension is determined by the generated
tracial von Neumann algebra.
\end{corollary}

\begin{proof}
For a bounded self-adjoint $d$-tuple $Z$, one has
$\delta(Z)\leq d$ and $\delta_0(Z)\leq d$. Both dimensions equal
$n$ for a freely independent variance-one semicircular $n$-tuple;
see~\cite[Section~2.4(a), (c), and (d)]{Voiculescu2002Entropy}
for $\delta$ and~\cite[Section~5]{BrannanEtAl2023} for $\delta_0$.
The standard semicircular $n$-tuple generates $L(\mathbb F_n)$ as
a von Neumann algebra. A self-adjoint pair supplied by
Theorem~\ref{thm:main} generates the same factor and has each
dimension at most $2<n$. These two tuples give the assertion.
\end{proof}

A separate isomorphism argument in
\emph{An isomorphism of the free group factors}
\cite[Corollary~7.1]{OpenAIFreeGroupFactors2026} gives a stronger
fixed-factor conclusion: each of $\delta,\delta_0$ and the
nonmicrostates dimensions $\delta^*,\delta^\star$ attains every
integer value $n\geq2$ on finite self-adjoint $W^*$-generating tuples
of $L(\mathbb F_2)$. That conclusion uses the trace-preserving
isomorphisms between free group factors established in the cited
paper. Corollary~\ref{cor:entropy-generators} uses only the
single-generation theorem proved here and the standard dimension
bounds above.
